% Generated draft: requires mathematical and source review.
\documentclass[11pt]{article}
\usepackage[margin=25mm]{geometry}
\usepackage{fontspec}
\setmainfont{Latin Modern Roman}
\usepackage{amsmath,amssymb,mathtools,unicode-math}
\setmathfont{Latin Modern Math}
\usepackage{enumitem,needspace,xurl,hyperref}
\setlength{\parindent}{0pt}
\setlength{\parskip}{5pt}
\setlength{\emergencystretch}{3em}
\newcommand{\N}{\mathbb{N}}
\newcommand{\Z}{\mathbb{Z}}
\newcommand{\Q}{\mathbb{Q}}
\newcommand{\R}{\mathbb{R}}
\newcommand{\C}{\mathbb{C}}
\newcommand{\F}{\mathbb{F}}
\newcommand{\A}{\mathbb{A}}
\newcommand{\PP}{\mathbb{P}}
\newcommand{\E}{\mathbb{E}}
\newcommand{\eps}{\varepsilon}
\newcommand{\dd}{\,\mathrm{d}}
\begin{document}
\section*{UnsolvedMath ID: 20001987}
\textbf{Source problem code:} catalogue entry
\subsection{Definitions and mathematical statement}
Let \(R_d=\Z[x_1^{\pm1},\ldots,x_d^{\pm1}]\), and let \(\mathfrak p\subset R_d\) be prime with \(\mathfrak p\cap\Z=0\). Write \(r=\dim(R_d/\mathfrak p)\) for the arithmetic Krull dimension and \(X=V(\mathfrak p\otimes\Q)\subset(\mathbb G_m)^d_\Q\) for the generic fiber, of dimension \(r-1\).

Fix the valuation-sign convention \(\operatorname{Log}_v(x)=(-\log|x_1|_v,\ldots,-\log|x_d|_v)\), with \(|q|_q=q^{-1}\). Define \(\mathcal A_v(X)\) as the closure of these images over \(\C\) at the archimedean place and over \(\overline\Q_q\) at a rational prime \(q\), and put \(\mathcal A_{\mathrm{ad}}(X)=\bigcup_v\mathcal A_v(X)\). With the opposite logarithmic sign all sets and halfspaces must be reflected. The source does not specify this convention or nonarchimedean closure.

Einsiedler and Lind ask whether every halfspace of every \((d-r+1)\)-dimensional linear subspace of \(\R^d\) meets the adelic amoeba. For the interpretation addressed by the supplied aid, a halfspace is \(\{u\in V:\lambda(u)>0\}\), with \(0\ne\lambda\in V^*\). This is an open halfspace with boundary through the origin; a closed halfspace or a ray including its origin is different. The source's principal-ideal ray remark leaves that endpoint convention unstated.

The aid attributes a counterexample to Payne, Example 1.3, and reports the following family. In \(R_4\), for any integer \(m\notin\{0,1\}\), set
\[
\mathfrak p_m=(x_2-x_1+1,\ x_3-x_1+m).
\]
The quotient identifies with
\[
\Z[z^{\pm1},(z-1)^{-1},(z-m)^{-1},w^{\pm1}],
\]
so the reported ideal is prime, has contraction zero to \(\Z\), and has Krull dimension three. Its generic fiber has dimension two, and the complementary subspace in the question therefore also has dimension two. With the valuation-sign convention above, the reported missed open halfspace is
\[
H=\R_{>0}(e_1+e_2+e_3)+\R e_4,
\qquad H\cap\mathcal A_{\mathrm{ad}}(X_m)=\varnothing.
\]
The \(m=2\) instance is the cited Payne example. The aid further reports that for every rational prime \(q\nmid m(m-1)\), the closed nonarchimedean amoeba contains positive coordinate rays in all four coordinate directions. Thus these amoebas are not contained in a linear hyperplane. Ray closure is material because coordinate logarithms of algebraic nonarchimedean points lie in rational multiples of the logarithm of the prime before taking closure; that rational value group is dense, not discrete.

The conclusion refutes this open-halfspace interpretation in the exact arithmetic dimension convention. It is not a counterexample to a differently oriented logarithmic convention without reflecting \(H\), nor automatically to the closed-halfspace version. The source's principal-ideal observation and the claimed generalization must be read with these distinctions.
\subsection{Short English statement}
Must adelic amoebas meet every complementary-dimensional halfspace through the origin? The supplied aid reports prime integral counterexamples in four variables, including Payne's example. The refutation uses open halfspaces and a specified valuation sign; arithmetic and generic-fiber dimensions differ by one.
\subsection{Sources}
\begin{itemize}
\item Manfred Einsiedler and Doug Lind, A.5, \emph{Half-space behavior of amoebas}, AIM workshop \emph{Amoebas and tropical geometry}, version dated 14 January 2004: \url{https://aimath.org/WWN/amoebas/amoebas.pdf}.
\item Payne, number-field Example 1.3, as named by the supplied aid; no title, year or identifier is supplied.
\item Local registry \texttt{data/unsolvedmath/problems.json}, permanent ID 20001987, code \texttt{catalogue entry}; canonical \texttt{aim-geometry-notes.json}, \texttt{notes[324]}; updated 14 August 2026. Clean statement null; statement status reconstructed unverified. The source tag section contains a genuine question; the original literature field is empty.
\item Attempt 1: counterexample, problem status at run refuted, rigor label proved, novelty confidence low; current solved and solved-in-literature labels. The aid is machine generated, not peer reviewed, and requires independent verification before citation.
\end{itemize}
\end{document}
